\section{Capítulo~\ref{sec:livingmachine}. Una máquina de neuronas vivas}

El comportamiento de los cultivos sobre matrices de electrodos está medido en experimentos directos con registro y estimulación; el mecanismo de las ráfagas se apoya en el modelo de agotamiento sináptico y en experimentos con microcultivos, y las afirmaciones sobre la dopamina en la placa, en experimentos aislados, algunos de los cuales los propios autores consideran solo una demostración.

{\footnotesize
\setlength{\tabcolsep}{3pt}\renewcommand{\arraystretch}{1.15}
\begin{longtable}{>{\raggedright}p{0.5cm}>{\raggedright}p{4.0cm}>{\raggedright}p{5.6cm}>{\raggedright}p{2.3cm}>{\raggedright\arraybackslash}p{1.7cm}}
\toprule
N.º & Afirmación & Experimento y qué se midió & Fuente & Estado \\
\midrule\endfirsthead
\toprule
N.º & Afirmación & Experimento y qué se midió & Fuente & Estado \\
\midrule\endhead
1 & Cultivo sobre chip: sobre todo \nt{GLUT}, con una fracción menor de \nt{GABA} que depende de la preparación; en cultivos de hipocampo, alrededor del $6\%$ & Una red de miles de neuronas sobre una matriz de electrodos es un experimento estándar (fila~3). La fracción de neuronas \nt{GABA} en cultivos de hipocampo se contó por tinción de GABA: alrededor del $6\%$ de las neuronas. Para cultivos de corteza no damos una cifra verificada & \cite{Benson1994} & medido \\
2 & Organoides: tejido nervioso tridimensional de alrededor de un milímetro, a partir de células madre humanas & A partir de células madre se cultivaron organoides tridimensionales con varias regiones cerebrales, entre ellas corteza con capas de progenitores y neuronas maduras & \cite{Lancaster2013} & medido \\
3 & Sin entrada, el cultivo se enciende en ráfagas con intervalos de un segundo a minutos, según el cultivo & Se registraron durante cinco semanas $58$ cultivos de corteza con densidades de $3\,000$ a $50\,000$ neuronas ($963$ registros de media hora): tras dos semanas, en casi todos los cultivos predominan las ráfagas de toda la red; los intervalos entre ráfagas van de $1$ a $300$~s y dependen mucho de la preparación & \cite{Wagenaar2006} & medido \\
4 & La ráfaga termina por agotamiento del neurotransmisor, con intervalo $T_{\text{ráfaga}}\approx\tau_{\text{rec}}\ln\frac{1}{1-x^*}$~\eqref{eq:burstinterval}; $2$--$4$~s es la estimación del modelo con $\tau_{\text{rec}}=0.8$~s, la recuperación medida es más lenta & La fórmula se deduce en el capítulo. Modelo: una red con sinapsis que se agotan genera por sí sola ráfagas de toda la red. Experimento en microcultivos de $4$--$30$ neuronas: la duración de la ráfaga depende del tiempo transcurrido desde la anterior, y el agotamiento de las vesículas de neurotransmisor suprime las ráfagas; conclusión de los autores: la ráfaga está limitada por la reserva de neurotransmisor. Allí la recuperación de la reserva tarda decenas de segundos, lo que con la misma fórmula da intervalos de decenas de segundos y minutos & deducción en el capítulo; \cite{Tsodyks2000,CohenSegal2011,Tsodyks1997} & teoría \\
5 & “Un maestro que se calla”: la estimulación se apaga cuando aparece la respuesta deseada, y la respuesta se fija & Se estimuló el cultivo a $0.3$--$1$~Hz hasta que aparecía la respuesta elegida $50\pm10$~ms después del estímulo; luego se apagaba la estimulación durante $5$~min; tras varios ciclos la respuesta deseada se provocaba enseguida; se trazaron curvas de aprendizaje & \cite{Shahaf2001} & medido \\
6 & El cultivo juega al tenis; un aumento significativo de las series en la sesión, solo con retroalimentación completa & En detalle, en el capítulo~\ref{sec:dishbrain} y su apéndice: con silencio en lugar del estímulo, el cultivo al final de la sesión también es mejor que sin retroalimentación, pero con un efecto menor; el aprendizaje de una sesión a otra es inestable & \cite{Kagan2022} & medido \\
7 & Que el cultivo aprenda a predecir su entrada es una hipótesis; las palabras grandilocuentes sobre su “sintiencia” están cuestionadas & El principio de energía libre es una propuesta teórica; no hay un experimento directo que lo distinga de explicaciones más sencillas. Treinta investigadores señalaron en un artículo de respuesta que un cambio de comportamiento en el lazo no muestra ni inteligencia ni sensaciones & \cite{Friston2010,Balci2023} & hipótesis \\
8 & El cultivo conduce un cuerpo virtual hacia un objetivo con un patrón de estímulos elegido & Un programa elegía por sí mismo los patrones de estimulación de entrenamiento según el resultado actual; en decenas de minutos la red alcanzaba los estados deseados, y la plasticidad se mantenía $80$~min tras $2$~h de entrenamiento. Los mismos estímulos, aplicados sin relación con el comportamiento, no tenían efecto & \cite{Bakkum2008} & medido \\
9 & Reservorio: solo se entrena la lectura lineal $y=W_{\text{sal}}\mathbf r$~\eqref{eq:reservoir} & Teoría de la computación de reservorio: cualquier sistema no lineal con memoria que se desvanece más una salida lineal entrenada & \cite{Maass2002,JaegerHaas2004} & teoría \\
10 & Un organoide como reservorio distingue hablantes con una exactitud de alrededor del $78\%$ (según los autores); por el error aprende la lectura, y las conexiones del organoide cambian sin maestro & Sonidos del habla de varios hablantes se entregaron al organoide como patrones de corriente en la matriz de electrodos; la lectura entrenada distinguía al hablante con una exactitud de alrededor del $78\%$, según informan los autores. Ellos mismos informan que durante el entrenamiento cambiaba la conectividad funcional del organoide, y lo llaman aprendizaje no supervisado en el propio organoide & \cite{Cai2023} & medido \\
11 & Un millón de neuronas gasta en espigas unos $10^{-4}$~W~\eqref{eq:dishpower} & Estimación: el coste de una espiga, $10^{-10}$~J, del presupuesto energético de la corteza~\eqref{eq:energy}, multiplicado por el número de espigas; no hay medición directa de la potencia del cultivo & deducción en el capítulo; \cite{Attwell2001} & teoría \\
12 & Dopamina con un destello de luz: $1$~mM de dopamina enjaulada, $240$ repeticiones, el número de espigas evocadas crecía & En una plataforma remota con organoides, la dopamina en una jaula fotosensible ($1$~mM) se liberaba con ultravioleta ($365$~nm, $0.8$~s) cuando el estímulo provocaba más espigas; a lo largo de $240$ pasos (unas $5$~h) el número de espigas creció en conjunto. Los autores escriben que con un solo experimento no se puede establecer la relación con la dopamina; no hay controles & \cite{Jordan2024} & hipótesis \\
13 & La dopamina de células vivas cambia el peso de un transistor orgánico de forma duradera y reversible & Células que liberan dopamina se cultivaron sobre un elemento neuromórfico orgánico; la oxidación de la dopamina en el electrodo cambia la conductancia del canal; se mostraron un cambio duradero del peso y su retorno & \cite{Keene2020} & medido \\
14 & Existen organoides con neuronas \nt{DOPA} funcionales; su dopamina no se ha usado como maestro & En organoides de células madre humanas se encontraron neuronas \nt{DOPA} maduras eléctricamente activas y producción de dopamina. No hemos encontrado en la literatura experimentos en los que esta dopamina enseñe a otra red & \cite{Jo2016} & medido \\
15 & Un organoide equilibra un péndulo: los estímulos de enseñanza del programa son mejores que los aleatorios, sin consolidación, a través de sinapsis \nt{GLUT} & Organoides de corteza de ratón en lazo cerrado con la tarea del “péndulo sobre un carro”: en la mayoría de los organoides, las señales elegidas por aprendizaje por refuerzo dieron mejor resultado que las aleatorias o que ninguna; tras $45$~min de descanso la mejora no se conservaba; el bloqueo de los receptores AMPA y NMDA la anulaba & \cite{Robbins2026} & medido \\
16 & Sin registros de control, la red del banco no puede decidir por sí misma qué cuenta como éxito (proposición~\ref{prop:livingmachine}) & En todos los experimentos de las filas 5--15 la señal de enseñanza la fija el experimentador o un programa; no hay ningún experimento en el que el cultivo haya elaborado por sí mismo un criterio de éxito: es una conclusión por ausencia, no un experimento & --- & hipótesis \\
\bottomrule
\end{longtable}}
